学会凸集和凸函数后，下一步要识别一个现实问题能否通过等价表述落入凸优化，模型名称将在这个过程中出现。许多问题表面上有范数、分式、最大值或矩阵不等式；只要保存其结构，便能转化为成熟算法可以处理的标准形式。

本单元从最透明的模型逐步扩展：

\begin{enumerate}
\tightlist
\item
  标准形式与线性规划说明为什么凸不等式必须朝正确方向，仿射等式为何特殊，并从多面体几何理解线性规划。
\item
  最小二乘、范数与二次规划把拟合误差、最坏偏差和二次代价放在同一条建模主线上。
\item
  变量变换揭示隐藏凸性以几何规划为例，展示一个原变量中看似非凸的问题怎样在对数坐标中变成凸问题。
\item
  半正定规划优化矩阵不等式用线性矩阵不等式表达谱、协方差和二次型约束，并说明凸松弛的力量与边界。
\end{enumerate}

标准形式是现实问题与求解算法之间的中间语言，无须要求原问题一开始就呈现这种外观。建模的关键始终是：变换是否等价、定义域是否保存，以及最终得到的凸性是否有严格依据。

\subsection{四类问题用一个模型互相对照}

为感受不同建模路径，考虑同一个数据拟合骨架。给定 \(A=\begin{bmatrix}1&2\\3&4\\5&6\end{bmatrix}\)，\(y=(4,10,16)^{\trans}\)。最小化二乘残差 \(\norm{Ax-y}_2^2\) 得 \(x^\star=(2,1)^{\trans}\)。若改用 \(L_1\) 范数 \(\norm{Ax-y}_1\)，目标不可微但凸，最优残差在坐标零附近产生稀疏模式。若要求 \(x_2\ge0\)，加入线性不等式后仍为凸规划。若进一步要求 \(x\) 是整数，可解性消失，问题进入整数规划领域。四类约束在同一个小例中连续分化，正是本单元结构的缩影。

\subsection{怎么读这一章}

这一章的重点不是看见一个漂亮公式就宣布它凸，而是依次检查变量的定义域、目标函数、每个不等式和每个等式。前两篇从常见标准模型入手，后两篇展示凸性有时藏在变量表示或矩阵不等式后面。需要补函数判据时，可回看 \hyperref[ch:073]{``凸函数：局部信息为何能控制全局''}。

\subsection{本章篇目}

以下篇目按概念依赖排列。

\begin{itemize}
\tightlist
\item \hyperref[sec:09-Convex-Optimization/04-Convex-Problems/01]{``标准形式与线性规划''}；
\item \hyperref[sec:09-Convex-Optimization/04-Convex-Problems/02]{``最小二乘、范数与二次规划''}；
\item \hyperref[sec:09-Convex-Optimization/04-Convex-Problems/03]{``变量变换揭示隐藏凸性''}；
\item \hyperref[sec:09-Convex-Optimization/04-Convex-Problems/04]{``半正定规划优化矩阵不等式''}。
\end{itemize}
